\section{Reconciled computations and negative optimization result}\label{sec:experiments}
\subsection{Evidence selection and elementary controls}
The repository retains both successful and failed campaigns. We reconstructed tables from recorded configurations, per-run outputs and training count traces, rather than from the previous manuscript. A source-hash inventory distinguishes byte-verified outputs, source-snapshot gaps, preserved failed attempts and unattributed copies. Historical driver scripts that overwrite fixed paths were not rerun. The new reconciliation is a deterministic analysis of the archive, not an independent repetition of every optimization trajectory.

An independent exhaustive check of all $64$ simple unweighted four-vertex graphs verifies costs, complementary partitions and exact minima. \Cref{tab:graphs} records representative graphs. Haar averaging gives a uniform law on the $D=\binom{m+k-1}{k}$ Fock compositions, so $\E P_{\CF}=\binom{m}{k}/D$ and $\E\succp=g^*/D$, where $g^*$ is the number of optimal valid \emph{strings}. At $m=4,k=2$, $D=10$. All six balanced strings are optimal on $K_4$, giving $\E\succp=0.6$; $K_{2,2}$ has minimum cut $2$ and four optimal strings, giving $0.4$. A graph-independent Haar success value $0.2$ is therefore incorrect. The Haar baseline is an ensemble average, not a universal per-circuit ceiling~\cite{arkhipov2012}.
\begin{table}[t]
\centering
\begin{tabular}{lrrc}
\hline
Graph & $C^*$ & Optimal strings $g$ & Haar $\mathbb{E}[P_{\rm opt}]$ \\
\hline
Path & 1 & 2 & $\frac{2}{10}$ \\
Cycle & 2 & 4 & $\frac{4}{10}$ \\
$K_4$ & 4 & 6 & $\frac{6}{10}$ \\
$K_{2,2}$ & 2 & 4 & $\frac{4}{10}$ \\
Star & 2 & 6 & $\frac{6}{10}$ \\
\hline
\end{tabular}
\caption{Computed four-vertex, two-photon controls. Balanced bitstrings and their complements are counted separately. The last column is a Haar ensemble expectation, not a trained or measured hardware success rate.}
\label{tab:graphs}
\end{table}


The E01 Haar archive includes size-dependent circuit counts; in particular the $n=10$ sample is smaller than the $n\le8$ samples. It does not provide trained $n=10,12$ optimization evidence. E02 uses the optimal-state mask during optimization and is retained as an oracle-assisted state-preparation diagnostic. E03--E04 compare small-instance objective/optimizer variants, while E06--E09 test structural baselines, spectral information, coherence diagnostics and label sensitivity. None warrants extrapolating a favorable oracle/Haar ratio to an optimum-blind algorithm.

\subsection{Exactly recorded query and attempt budgets}
The primary empirical comparison is E05b. The full grid has $n=4,6,8$, six graph families, two nominal graph seeds, three training seeds and seven method configurations: bosonic and distinguishable meshes, a classical iid multinomial control, a uniform feasible sampler, and classical product-Bernoulli penalty controls at three strengths ($0.1,1,10$). Repeated graph seeds can generate identical graphs; the nominal twelve graph positions at each size contain only $9,9,10$ unique edge lists, respectively. The archived global means retain the nominal grid weighting so the analysis does not select favorable graphs.

Each optimized trajectory records $120$ objective evaluations, $256$ ideal sampled attempts per evaluation, and a final $4096$-attempt batch: $120\times256+4096=34816$ attempts. This equality does not make classical matrix simulation costs equal to physical circuit shots or to a discrete-search query. Mesh dimension is obtained from its actual coupler list and phases, not equated to the number of Bernoulli parameters. Rejected/colliding outcomes remain in the denominator. The retrospective optimum is used for evaluation, not the E05b recorded training loss; anti-oracle tests enforce that separation in the recovered public training API.

\begin{table}[t]
\centering\scriptsize
\setlength{\tabcolsep}{3pt}
\begin{tabular}{lccc}
\hline
Method & $n=4$ & $n=6$ & $n=8$ \\
\hline
Bosonic & \shortstack{0.652\\{[0.464, 0.804]}} & \shortstack{0.280\\{[0.178, 0.423]}} & \shortstack{0.068\\{[0.033, 0.107]}} \\
Distinguishable & \shortstack{0.713\\{[0.590, 0.835]}} & \shortstack{0.537\\{[0.389, 0.646]}} & \shortstack{0.138\\{[0.075, 0.197]}} \\
Classical iid multinomial & \shortstack{0.421\\{[0.310, 0.532]}} & \shortstack{0.110\\{[0.063, 0.158]}} & \shortstack{0.026\\{[0.016, 0.036]}} \\
Product Bernoulli penalty ($\lambda=0.1$) & \shortstack{0.199\\{[0.135, 0.268]}} & \shortstack{0.057\\{[0.033, 0.089]}} & \shortstack{0.017\\{[0.010, 0.023]}} \\
Product Bernoulli penalty ($\lambda=1$) & \shortstack{0.203\\{[0.141, 0.271]}} & \shortstack{0.056\\{[0.032, 0.088]}} & \shortstack{0.016\\{[0.010, 0.023]}} \\
Product Bernoulli penalty ($\lambda=10$) & \shortstack{0.199\\{[0.127, 0.270]}} & \shortstack{0.056\\{[0.032, 0.092]}} & \shortstack{0.017\\{[0.010, 0.023]}} \\
Uniform feasible & \shortstack{0.528\\{[0.333, 0.718]}} & \shortstack{0.175\\{[0.100, 0.277]}} & \shortstack{0.060\\{[0.036, 0.086]}} \\
\hline
\end{tabular}
\caption{Ideal final success means over 36 runs per method and size, with descriptive graph-cluster percentile intervals. There are 12 nominal graph slots per size but $G=9,9,10$ unique labeled weighted graphs, respectively. Every run uses 120 calls of 256 training draws plus 4096 heldout draws (34816 ideal attempts). Optimizer repeats stay in the same graph cluster; these are not population significance intervals. Each penalty parameter is a separate fully charged run. No hardware shots.}
\label{tab:e05b}
\end{table}

\begin{figure}[tb]
 \centering\includegraphics[width=.94\linewidth]{figures/fig_e05b.pdf}
 \caption{Final success of bosonic and distinguishable E05b procedures, with the complete nominal-grid means shown as points. Descriptive intervals resample unique fixed-graph units after averaging training seeds; they are not a population confidence claim for random graph ensembles. Per-graph and per-seed values, penalty settings and aggregation rules are supplied with the data.}\label{fig:e05b}
\end{figure}

The bosonic final means are lower than the distinguishable means at all three sizes. These observations reject the claimed benefit for this implementation on this grid and budget; they do not prove that every optical optimizer is inferior. Every archived run already encounters an optimal solution during training, so the final circuit probability alone fails to describe problem-solving utility. Count traces locate the first \emph{winning batch}, but do not preserve shot order inside a batch. Exact first-hit shot counts cannot be recovered and are not imputed. Thus a claim of shorter time to first optimum is not supported by these data. A future search trial must save ordered queries and charge preprocessing, restarts, training, final samples and decoding.

The uniform balanced sampler is a relevant feasibility baseline; Bernoulli penalties are a different, weaker control and their three strengths are reported separately. The P21 classical search archive supplies balanced random search, exchange/local improvement and graph-aware controls at a separate cut-query budget. Spectral preprocessing and graph relabeling tests show why a successful optically encoded sample can inherit classically supplied graph structure. These campaigns strengthen the no-benefit assessment, but their differing budgets do not form a matched wall-clock speedup experiment. The new evidence inventory gives their exact provenance and resource units.

ObliQ is an actual hybrid photonic QUBO algorithm~\cite{ranjan2025}, not a synonym for an independent Bernoulli sampler. The authenticated static optical kernel was checked against its archive, Fock calculations and Perceval; literal and physically corrected angle conventions are kept separate. Complete original training, source selection and prefix decoding have not been reproduced, so no end-to-end superiority over authentic ObliQ is claimed.

\subsection{Finite theory checks and retained counterexamples}
The accepted supervised P33 primary/replication data cover all $300$ registered grid circuits, with full-law computations restricted to at most twelve modes; an additional $72$ entries evaluate certificate quantities only. Larger-frame entries evaluate certificate quantities only. Independent complete-law and Perceval controls support implementation consistency within stated tolerances, without proving the theorem. \Cref{fig:density} is generated from those accepted files; unattributed root-level P33 copies are excluded.
\begin{figure}[tb]
 \centering\includegraphics[width=\linewidth]{figures/fig_density.pdf}
 \caption{Finite P33 certificate and sharpness checks from the accepted supervised primary/replication archive. Full-law comparisons and certificate-only large frames are distinguished by the underlying record; no large-size full-Fock claim is inferred. The two-photon sharp family and dilute spectator obstruction test different hypotheses of \cref{thm:density}.}\label{fig:density}
\end{figure}

The P35 low-rank threshold expectation calculation is an exact classical mean-objective control, consistent with known product-state counting methods~\cite{ivanov2020}; it is not a sampler or a new optimization algorithm. Its earlier resource-guard defect and independent post-fix receipt are preserved. A stronger proposed pinched-state curved envelope was refuted in P39 by an exact physical three-photon/nine-mode construction at leakage $18/25$. Its measured output total variation remains below that curve: the refutation concerns the larger auxiliary state metric. The general measured bound $\TV\le\delta$ remains open. Failed searches and nonconverged solvers are retained as such, not promoted into proofs.

